% ============================================================
%  Rapport d'avancement — Benchmark IA multilangue
%  HACA / Direction des Systèmes d'Information
%  Auteur : Marwane ElBaraka
%  Date   : 16 juin 2026
%  Compiler with: pdflatex (twice) or upload to Overleaf
% ============================================================
\documentclass[12pt, a4paper]{article}

% ---------- Encodage & langue ----------
\usepackage[utf8]{inputenc}
\usepackage[T1]{fontenc}
\usepackage[french]{babel}

% ---------- Mise en page ----------
\usepackage[top=2.5cm, bottom=2.5cm, left=2.8cm, right=2.8cm]{geometry}
\usepackage{setspace}
\setstretch{1.15}
\usepackage{parskip}          % espace entre paragraphes, pas d'indentation
\setlength{\parindent}{0pt}

% ---------- Typographie ----------
\usepackage{lmodern}
\usepackage{microtype}

% ---------- Couleurs & boîtes ----------
\usepackage[dvipsnames, table]{xcolor}
\definecolor{hacablue}{RGB}{0, 70, 127}
\definecolor{lightgray}{RGB}{245, 245, 245}
\definecolor{darkgray}{RGB}{80, 80, 80}
\definecolor{goodgreen}{RGB}{0, 150, 80}
\definecolor{warnorange}{RGB}{220, 100, 0}

% ---------- Tableaux ----------
\usepackage{booktabs}
\usepackage{array}
\usepackage{tabularx}
\usepackage{multirow}
\usepackage{colortbl}

% ---------- Graphiques ----------
\usepackage{graphicx}
\usepackage{float}
\usepackage{caption}
\captionsetup{font=small, labelfont=bf, skip=6pt}

% ---------- Listes ----------
\usepackage{enumitem}
\setlist[itemize]{noitemsep, topsep=3pt}
\setlist[enumerate]{noitemsep, topsep=3pt}

% ---------- Code ----------
\usepackage{listings}
\lstset{
  basicstyle=\ttfamily\footnotesize,
  backgroundcolor=\color{lightgray},
  frame=single, framerule=0.4pt,
  breaklines=true, columns=flexible
}

% ---------- Liens ----------
\usepackage[hidelinks, colorlinks=true,
            linkcolor=hacablue, urlcolor=hacablue]{hyperref}

% ---------- En-tête / pied de page ----------
\usepackage{fancyhdr}
\pagestyle{fancy}
\fancyhf{}
\fancyhead[L]{\small\textcolor{darkgray}{Rapport d'avancement — Benchmark IA multilangue}}
\fancyhead[R]{\small\textcolor{darkgray}{HACA / DSI — \the\year}}
\fancyfoot[C]{\small\thepage}
\renewcommand{\headrulewidth}{0.4pt}

% ---------- Titre section coloré ----------
\usepackage{titlesec}
\titleformat{\section}
  {\large\bfseries\color{hacablue}}{\thesection.}{0.6em}{}[\vspace{-2pt}\rule{\linewidth}{0.5pt}]
\titleformat{\subsection}
  {\normalsize\bfseries\color{hacablue!80!black}}{\thesubsection.}{0.5em}{}
\titleformat{\subsubsection}
  {\normalsize\itshape\color{darkgray}}{\thesubsubsection.}{0.4em}{}

% ---------- Boîtes info ----------
\usepackage{tcolorbox}
\tcbuselibrary{skins, breakable}
\newtcolorbox{infobox}[1][]{
  colback=hacablue!5, colframe=hacablue!60,
  fonttitle=\bfseries\small, title=#1,
  boxrule=0.5pt, arc=3pt
}
\newtcolorbox{resultbox}{
  colback=goodgreen!7, colframe=goodgreen!50,
  boxrule=0.5pt, arc=3pt
}

% ============================================================
\begin{document}

% -------- PAGE DE GARDE --------
\begin{titlepage}
  \centering
  \vspace*{1.5cm}

  {\LARGE\bfseries\color{hacablue}
    Rapport d'avancement\\[0.4em]
    \large Benchmark de modèles d'IA pour l'analyse de tonalité\\
    dans les médias audiovisuels marocains
  }\par

  \vspace{1.2cm}
  \rule{0.6\linewidth}{1pt}\par
  \vspace{0.6cm}

  {\large
    \textbf{Stage d'initiation} — Direction des Systèmes d'Information\\[0.3em]
    \textbf{Organisation :} Haute Autorité de la Communication Audiovisuelle (HACA)\\[0.3em]
    \textbf{Stagiaire :} Marwane ElBaraka\\[0.3em]
    \textbf{Date :} 16 juin 2026
  }\par

  \vspace{1.2cm}
  \rule{0.6\linewidth}{1pt}\par
  \vspace{1.5cm}

  \begin{infobox}[Résumé exécutif]
    Ce rapport présente l'avancement d'un benchmark comparatif de sept modèles
    de traitement automatique du langage naturel (TALN) pour la classification
    de tonalité (positif / neutre / négatif) sur des contenus audiovisuels en
    quatre langues : darija marocain (écriture arabe), arabizi marocain (écriture
    latine), arabe standard moderne et français. \textbf{Quatre étapes sur sept
    sont complètes.} Les modèles fine-tunés dépassent les modèles génériques sur
    toutes les langues cibles, avec un macro-F1 maximal de \textbf{0{,}983} pour
    l'arabizi et \textbf{0{,}844} pour le darija.
  \end{infobox}

  \vfill
  {\small\textcolor{darkgray}{Document confidentiel — usage interne HACA}}
\end{titlepage}

% -------- TABLE DES MATIÈRES --------
\tableofcontents
\newpage

% ============================================================
\section{Contexte et objectifs}
% ============================================================

\subsection{Présentation de la HACA}

La Haute Autorité de la Communication Audiovisuelle (HACA) est l'autorité
marocaine de régulation du secteur audiovisuel, créée en 2002. Elle est
chargée de veiller au respect du pluralisme, de la déontologie et des
obligations légales des opérateurs de radio et de télévision, ainsi
qu'accompagner la transformation numérique du secteur.

La Direction des Systèmes d'Information (DSI) est l'entité responsable
du développement et de la maintenance des outils informatiques internes.
Elle développe notamment des plateformes d'échange avec les opérateurs
audiovisuels (\textit{HACA Bridges}) et des solutions logicielles
commercialisées auprès d'autres régulateurs africains (\textit{HMS —
HACA Media Solutions}).

\subsection{Problématique}

Le suivi de la tonalité éditoriale des contenus diffusés constitue une
mission clé de toute autorité de régulation audiovisuelle. Les contenus
marocains sont produits dans au moins \textbf{quatre variétés linguistiques} :

\begin{itemize}
  \item \textbf{Arabe standard moderne (MSA)} — journaux télévisés, débats officiels
  \item \textbf{Darija marocain en écriture arabe} — programmes de divertissement, réseaux sociaux
  \item \textbf{Arabizi marocain (écriture latine)} — commentaires YouTube, forums en ligne
  \item \textbf{Français} — médias francophones, journaux télévisés bilingues
\end{itemize}

Aucun outil existant ne couvre simultanément ces quatre variétés. L'objectif
du stage est de \textbf{construire un benchmark reproductible} permettant
d'identifier les meilleurs modèles open-source pour chaque langue, en vue
d'une intégration dans les outils de surveillance de la DSI.

\subsection{Métrique principale}

Le \textbf{macro-F1} est utilisé comme métrique principale : il calcule la
moyenne des F1 par classe, sans pondération par fréquence. Cette métrique
pénalise les modèles qui ignorent les classes minoritaires, ce qui est
essentiel dans un contexte où la classe « neutre » est souvent sous-représentée
mais reste informationnellement importante.

% ============================================================
\section{Architecture du benchmark}
% ============================================================

\subsection{Vue d'ensemble du projet}

Le projet est structuré en \textbf{7 étapes} séquentielles, dont
\textbf{4 sont complètes} à la date de ce rapport :

\begin{center}
\begin{tabular}{clll}
\toprule
\textbf{Étape} & \textbf{Contenu} & \textbf{Statut} \\
\midrule
1 & Mise en place du dépôt, environnement, tests de fumée & \textcolor{goodgreen}{\textbf{Terminé}} \\
2 & Construction et gel des jeux de test (4 langues) & \textcolor{goodgreen}{\textbf{Terminé}} \\
3 & Évaluation des 3 modèles prêts à l'emploi & \textcolor{goodgreen}{\textbf{Terminé}} \\
4 & Fine-tuning de 4 modèles spécialisés & \textcolor{goodgreen}{\textbf{Terminé}} \\
5 & Atlas-Chat (LLM zéro-shot, accès conditionnel) & \textcolor{warnorange}{\textit{En attente}} \\
6 & Corpus domaine réel (fichiers SRT HACA) & \textcolor{warnorange}{\textit{En attente}} \\
7 & Visualisations finales, grille pondérée, rapport & \textcolor{warnorange}{\textit{Partiel}} \\
\bottomrule
\end{tabular}
\end{center}

\subsection{Principes méthodologiques}

\begin{itemize}
  \item \textbf{Reproductibilité} : graine aléatoire \texttt{seed=42} partout
    (Python, NumPy, PyTorch, CUDA). Dépendances gelées dans \texttt{requirements.txt}.
  \item \textbf{Jeux de test immuables} : construits une seule fois, hashés en
    SHA-256, jamais modifiés par la suite.
  \item \textbf{Harness unifié} : même code d'évaluation pour tous les modèles
    (prêts à l'emploi ou fine-tunés), garantissant la comparabilité.
  \item \textbf{Évaluation par langue} : aucun score global agrégeant toutes les langues.
  \item \textbf{Mappage de labels centralisé} : vérification systématique de
    \texttt{model.config.id2label} avant toute interprétation des sorties.
\end{itemize}

% ============================================================
\section{Étapes réalisées}
% ============================================================

\subsection{Étape 1 — Mise en place de l'environnement}

Le dépôt Git a été initialisé avec la structure suivante :

\begin{lstlisting}
benchmark/
  src/           # modules Python (harness, finetune, utils, ...)
  data/
    raw/         # données brutes (gitignore)
    test_sets/   # jeux de test gelés (darija_ar.csv, francais.csv, ...)
  results/       # un JSON par (modele x langue) + summary.csv
  results/figs/  # graphiques générés
  checkpoints/   # poids des modèles fine-tunés
  report/        # documentation et rapports
  notebooks/     # notebook Kaggle (fine-tuning)
  Dockerfile     # reproducibilité totale
\end{lstlisting}

Un test de fumée (\textit{smoke test}) a été exécuté avec le modèle
\texttt{xlm-t} sur trois phrases de référence, validant l'environnement
d'inférence avant toute mesure formelle.

\subsection{Étape 2 — Construction des jeux de test}

\subsubsection{Sources de données}

Quatre corpus publics ont été utilisés pour constituer les jeux de test :

\begin{center}
\begin{tabularx}{\textwidth}{llXl}
\toprule
\textbf{Langue} & \textbf{Corpus} & \textbf{Source} & \textbf{Classes} \\
\midrule
darija\_ar & MAC & LeMGarouani/MAC (GitHub) & 3 (neg/neu/pos) \\
français & Allociné & HuggingFace datasets & 2 (neg/pos) \\
MSA & ASTD & mahmoudnabil/ASTD (GitHub) & 2 (neg/pos) \\
arabizi & MYC & MouadJb/MYC (GitHub), script latin uniquement & 2 (neg/pos) \\
\bottomrule
\end{tabularx}
\end{center}

\subsubsection{Protocole d'échantillonnage}

Pour chaque langue : tirage stratifié de \textbf{1\,000 lignes} (graine 42),
sauvegarde au format \texttt{[text, label]}, calcul du hash SHA-256 pour
garantir l'intégrité, vérification manuelle de la distribution des classes.

\begin{center}
\begin{tabular}{lrrrl}
\toprule
\textbf{Langue} & \textbf{neg} & \textbf{neu} & \textbf{pos} & \textbf{Remarque} \\
\midrule
darija\_ar & 201 & 232 & 567 & Déséquilibré, pos majoritaire \\
français   & 520 &  —  & 480 & Binaire, quasi-équilibré \\
MSA        & 679 &  —  & 321 & Binaire, neg majoritaire \\
arabizi    & 343 &  —  & 657 & Binaire, pos majoritaire \\
\bottomrule
\end{tabular}
\end{center}

\textbf{Remarque sur le corpus MYC (arabizi) :} seules les lignes à écriture
exclusivement latine (sans caractères arabes U+0600–U+06FF) ont été conservées,
afin d'isoler l'arabizi pur des autres variétés.

\subsection{Étape 3 — Modèles prêts à l'emploi}

Trois modèles pré-entraînés ont été évalués sans entraînement supplémentaire,
constituant la \textbf{ligne de base} du benchmark.

\subsubsection{xlm-t — Modèle multilingue généraliste}

\textbf{Modèle :} \texttt{cardiffnlp/twitter-xlm-roberta-base-sentiment}
(278M paramètres) — RoBERTa multilingue fine-tuné sur 198 millions de tweets
en 100+ langues. Seul modèle couvrant les quatre langues cibles sans
entraînement additionnel.

\begin{center}
\begin{tabular}{lrrr}
\toprule
\textbf{Langue} & \textbf{Macro-F1} & \textbf{Latence CPU (ms/utt)} & \textbf{Observation} \\
\midrule
darija\_ar & 0{,}723 & 109{,}6 & Classe neutre difficile (F1=0{,}577) \\
français   & 0{,}772 & 1006{,}1 & Latence anormale (textes longs) \\
MSA        & 0{,}813 & 168{,}4 & 12,9\,\% d'abstentions (classe neutre) \\
arabizi    & 0{,}762 & 145{,}3 & 31,2\,\% d'abstentions sur ce jeu binaire \\
\bottomrule
\end{tabular}
\end{center}

\textbf{Problème identifié :} xlm-t prédit « neutre » pour 31\,\% des
entrées arabizi, car ce modèle n'a pas été exposé à l'arabizi marocain
pendant son entraînement. Ce problème d'abstention est la principale
motivation pour les fine-tunings de l'étape 4.

\subsubsection{camelbert-da — Spécialiste arabe dialectal}

\textbf{Modèle :} \texttt{CAMeL-Lab/bert-base-arabic-camelbert-da-sentiment}
(109M paramètres) — BERT pré-entraîné sur de l'arabe dialectal multi-pays.

\begin{center}
\begin{tabular}{lrrr}
\toprule
\textbf{Langue} & \textbf{Macro-F1} & \textbf{Latence CPU (ms/utt)} & \textbf{Observation} \\
\midrule
darija\_ar & 0{,}701 & 90{,}9 & Biais vers le négatif (précision=0,581) \\
MSA        & \textbf{0{,}924} & 163{,}7 & Meilleur résultat MSA du benchmark \\
\bottomrule
\end{tabular}
\end{center}

\textbf{Observation notable :} ce modèle documente une tâche binaire mais
expose réellement trois classes (neg/neu/pos). Identifié via inspection de
\texttt{model.config.id2label} — cas typique de divergence entre documentation
et implémentation.

\subsubsection{distilcamembert — Spécialiste français}

\textbf{Modèle :} \texttt{cmarkea/distilcamembert-base-sentiment}
(68M paramètres) — DistilCamemBERT fine-tuné sur les critiques Allociné.

\begin{center}
\begin{tabular}{lrrr}
\toprule
\textbf{Langue} & \textbf{Macro-F1} & \textbf{Latence CPU (ms/utt)} & \textbf{Observation} \\
\midrule
français & \textbf{0{,}949} & 420{,}9 & Taux d'erreur de 2\,\% (20/1000) \\
\bottomrule
\end{tabular}
\end{center}

\textbf{Enseignement clé :} le plus petit modèle (68M) obtient le deuxième
meilleur score du benchmark, confirment que la spécialisation domine la taille.

\subsection{Étape 4 — Fine-tuning des modèles spécialisés}

\subsubsection{Protocole d'entraînement}

Tous les modèles ont été fine-tunés sur \textbf{Kaggle (GPU T4 16\,Go)}
avec la configuration suivante :

\begin{itemize}
  \item 3 époques, batch 16, learning rate $2 \times 10^{-5}$, fp16
  \item Longueur maximale : 128 tokens
  \item Sélection du meilleur checkpoint par macro-F1 sur validation
  \item Évaluation via le même harness que les modèles prêts à l'emploi
\end{itemize}

\subsubsection{Modèles fine-tunés}

\textbf{DarijaBERT} (\texttt{SI2M-Lab/DarijaBERT}, 147M paramètres) —
pré-entraîné spécifiquement sur le darija marocain (écriture arabe) par le
laboratoire SI2M (Université Mohammed V, Rabat). Fine-tuné sur MAC.
Durée d'entraînement : \textbf{11,6 min}.

\textbf{DarijaBERT-arabizi} (\texttt{SI2M-Lab/DarijaBERT-arabizi}, 171M paramètres) —
variante couvrant l'arabizi latin-marocain. Fine-tuné sur MYC-Latin.
Durée d'entraînement : \textbf{3,4 min}.

\textbf{MARBERTv2} (\texttt{UBC-NLP/MARBERTv2}, 163M paramètres) —
pré-entraîné sur 1 milliard de mots couvrant 25 dialectes arabes (dont le
marocain) + MSA. Fine-tuné sur MAC.
Durée d'entraînement : \textbf{11,8 min}.
\textcolor{warnorange}{\textbf{Licence : recherche uniquement (non-commerciale).}}

\textbf{QARIB} (\texttt{qarib/bert-base-qarib}, 135M paramètres) —
pré-entraîné sur 420M mots d'arabe du Golfe + MSA. Fine-tuné sur MAC.
Durée d'entraînement : \textbf{11,3 min}.

% ============================================================
\section{Résultats}
% ============================================================

\subsection{Tableau récapitulatif complet}

\begin{center}
\setlength{\tabcolsep}{5pt}
\begin{tabular}{lcccccc}
\toprule
\textbf{Modèle} & \textbf{darija\_ar} & \textbf{français} & \textbf{MSA} &
\textbf{arabizi} & \textbf{Params} & \textbf{Type} \\
\midrule
xlm-t           & 0{,}723 & 0{,}772 & 0{,}813 & 0{,}762 & 278M & Multilingue \\
camelbert-da    & 0{,}701 &  —      & \cellcolor{goodgreen!30}\textbf{0{,}924} &  —  & 109M & Arabe dial. \\
distilcamembert &  —      & \cellcolor{goodgreen!30}\textbf{0{,}949} &  —  &  —  & 68M  & Français \\
\midrule
DarijaBERT      & 0{,}775 &  —      &  —      &  —      & 147M & Fine-tuné \\
DarijaBERT-arabizi &  —   &  —      &  —      & \cellcolor{goodgreen!30}\textbf{0{,}983} & 171M & Fine-tuné \\
MARBERTv2       & \cellcolor{goodgreen!30}\textbf{0{,}844} &  — & 0{,}838 &  — & 163M & Fine-tuné \\
QARIB           & 0{,}827 &  —      & 0{,}812 &  —      & 135M & Fine-tuné \\
\bottomrule
\multicolumn{7}{l}{\small\textit{Cellules vertes = meilleur résultat par langue. Ligne supérieure = modèles prêts à l'emploi.}}
\end{tabular}
\end{center}

\subsection{Résultats par langue}

\subsubsection{Darija marocain (écriture arabe)}

\begin{resultbox}
\textbf{Meilleur modèle :} MARBERTv2 — Macro-F1 = \textbf{0{,}844}\\
\textbf{Amélioration vs ligne de base (xlm-t 0{,}723) :} +12{,}1 points
\end{resultbox}

La progression du F1 sur la classe neutre illustre l'impact de la spécialisation :

\begin{center}
\begin{tabular}{lrrrr}
\toprule
\textbf{Modèle} & \textbf{F1(neg)} & \textbf{F1(neu)} & \textbf{F1(pos)} & \textbf{Macro-F1} \\
\midrule
camelbert-da (base) & 0{,}708 & 0{,}526 & 0{,}868 & 0{,}701 \\
xlm-t (base)        & 0{,}724 & 0{,}577 & 0{,}868 & 0{,}723 \\
DarijaBERT (FT)     & 0{,}743 & 0{,}690 & 0{,}891 & 0{,}775 \\
QARIB (FT)          & 0{,}838 & 0{,}733 & 0{,}909 & 0{,}827 \\
MARBERTv2 (FT)      & \textbf{0{,}854} & \textbf{0{,}751} & \textbf{0{,}927} & \textbf{0{,}844} \\
\bottomrule
\end{tabular}
\end{center}

La classe neutre — la plus difficile — passe de 0{,}526 (camelbert-da) à
\textbf{0{,}751} (MARBERTv2), franchissant pour la première fois le seuil
de 0{,}75 fixé comme objectif dans le plan.

\subsubsection{Français}

\begin{resultbox}
\textbf{Meilleur modèle :} distilcamembert — Macro-F1 = \textbf{0{,}949}\\
\textbf{Taux d'erreur :} 2{,}0\,\% (20 erreurs sur 1\,000)
\end{resultbox}

Aucun fine-tuning n'a été réalisé pour le français à cette étape : le modèle
distilcamembert (prêt à l'emploi) constitue déjà une référence solide, avec
une avance de +17{,}7 points sur xlm-t. Un fine-tuning dédié (étape 7) pourrait
encore améliorer les performances sur du contenu audiovisuel réel.

\subsubsection{Arabe standard moderne (MSA)}

\begin{resultbox}
\textbf{Meilleur modèle :} camelbert-da (prêt à l'emploi) — Macro-F1 = \textbf{0{,}924}\\
\textbf{Taux d'erreur :} 3{,}2\,\% (32 erreurs sur 1\,000)
\end{resultbox}

Les modèles fine-tunés (MARBERTv2 0{,}838, QARIB 0{,}812) n'atteignent pas
camelbert-da sur le MSA. La raison est structurelle : ces modèles ont été
fine-tunés avec une tête à trois classes (neg/neu/pos) sur des données darija,
et lorsqu'ils sont évalués sur le jeu MSA binaire (neg/pos), leur classe
neutre absorbe 18--20\,\% des prédictions, réduisant le rappel. Un fine-tuning
dédié MSA avec une tête binaire serait la correction appropriée.

\subsubsection{Arabizi marocain}

\begin{resultbox}
\textbf{Meilleur modèle :} DarijaBERT-arabizi — Macro-F1 = \textbf{0{,}983}\\
\textbf{Taux d'erreur :} 1{,}5\,\% (15 erreurs sur 1\,000) — meilleur résultat du benchmark
\end{resultbox}

\begin{center}
\begin{tabular}{lrrrr}
\toprule
\textbf{Modèle} & \textbf{F1(neg)} & \textbf{F1(pos)} & \textbf{Macro-F1} & \textbf{Abstentions} \\
\midrule
xlm-t (base)            & 0{,}780 & 0{,}745 & 0{,}762 & 31{,}2\,\% \\
DarijaBERT-arabizi (FT) & \textbf{0{,}978} & \textbf{0{,}989} & \textbf{0{,}983} & 0\,\% \\
\bottomrule
\end{tabular}
\end{center}

L'amélioration de +22{,}1 points est la plus importante du benchmark. Elle
résout le problème des abstentions d'xlm-t (31\,\% d'entrées prédites
« neutre » par défaut) et produit le résultat le plus élevé, avec seulement
3,4 minutes d'entraînement sur GPU Kaggle T4.

\subsection{Latence et coût opérationnel}

Les modèles fine-tunés ont été mesurés sur GPU (Kaggle T4) ; les modèles
prêts à l'emploi sur CPU. Le gain est de l'ordre de \textbf{30 à 300×} :

\begin{center}
\begin{tabular}{llrr}
\toprule
\textbf{Catégorie} & \textbf{Matériel} & \textbf{Latence (ms/utt)} &
\textbf{Temps/émission 30 min} \\
\midrule
Modèles prêts à l'emploi & CPU & 90 -- 1\,006 & 5 -- 60 min \\
Modèles fine-tunés & GPU T4 & 2{,}5 -- 4{,}4 & 9 -- 16 s \\
\bottomrule
\end{tabular}
\end{center}

À raison de 2 sous-titres par seconde pour une émission de 30 minutes
(≈3\,600 utterances), les modèles fine-tunés permettent une analyse
quasiment en temps réel.

% ============================================================
\section{Visualisations produites}
% ============================================================

Trois graphiques ont été générés par le script \texttt{src/aggregate.py} :

\begin{enumerate}
  \item \textbf{Heatmap macro-F1} (\texttt{results/figs/heatmap\_f1.png}) —
    grille modèle × langue, couleur proportionnelle au score. Visualise d'un coup
    d'œil qu'aucun modèle ne domine toutes les langues, et que DarijaBERT-arabizi
    (0{,}983) et distilcamembert (0{,}949) sont les deux cellules les plus sombres.

  \item \textbf{Nuage coût-précision} (\texttt{results/figs/scatter\_cost\_precision.png}) —
    axe X : latence (log), axe Y : macro-F1 moyen, taille de bulle : racine carrée
    du nombre de paramètres. Montre clairement deux clusters : modèles GPU
    (2--5 ms, haute qualité) vs modèles CPU (90--1\,000 ms, qualité variable).

  \item \textbf{Radar des finalistes} (\texttt{results/figs/radar\_finalists.png}) —
    4 axes normalisés (Précision 50\,\%, Intégration 20\,\%, Coût 20\,\%, Couverture 10\,\%).
    Identifie DarijaBERT-arabizi comme premier au score pondéré (0{,}845),
    suivi de distilcamembert (0{,}724) et MARBERTv2 (0{,}540).
\end{enumerate}

\begin{center}
\begin{tabular}{lrrrrr}
\toprule
\textbf{Modèle} & \textbf{Précision} & \textbf{Coût} & \textbf{Couverture} &
\textbf{Score pondéré} & \textbf{Rang} \\
\midrule
DarijaBERT-arabizi & 1{,}000 & 0{,}599 & 0{,}25 & \textbf{0{,}845} & 1 \\
distilcamembert    & 0{,}839 & 0{,}400 & 0{,}25 & 0{,}724 & 2 \\
MARBERTv2          & 0{,}340 & 0{,}599 & 0{,}50 & 0{,}540 & 3 \\
camelbert-da       & 0{,}207 & 0{,}821 & 0{,}50 & 0{,}518 & 4 \\
QARIB              & 0{,}239 & 0{,}656 & 0{,}50 & 0{,}501 & 5 \\
xlm-t              & 0{,}000 & 0{,}491 & 1{,}00 & 0{,}398 & 6 \\
DarijaBERT         & 0{,}033 & 0{,}664 & 0{,}25 & 0{,}374 & 7 \\
\bottomrule
\multicolumn{6}{l}{\small\textit{Précision normalisée 50\,\%, Intégration 20\,\%, Coût 20\,\%, Couverture 10\,\%.}}
\end{tabular}
\end{center}

% ============================================================
\section{Analyses croisées et enseignements}
% ============================================================

\subsection{Le spécialiste bat le généraliste}

Le résultat le plus constant du benchmark est qu'un modèle spécialisé sur
une langue et un domaine surpasse systématiquement un modèle généraliste
plus grand :

\begin{center}
\begin{tabular}{lrrrr}
\toprule
\textbf{Langue} & \textbf{Généraliste} & \textbf{Spécialiste} & \textbf{Gain} \\
\midrule
Français     & xlm-t 0{,}772  & distilcamembert 0{,}949   & +0{,}177 \\
MSA          & xlm-t 0{,}813  & camelbert-da 0{,}924       & +0{,}111 \\
Darija AR    & xlm-t 0{,}723  & MARBERTv2 fine-tuné 0{,}844 & +0{,}121 \\
Arabizi      & xlm-t 0{,}762  & DarijaBERT-arabizi 0{,}983  & +0{,}221 \\
\bottomrule
\end{tabular}
\end{center}

\subsection{La classe neutre reste la plus difficile}

Sur le jeu darija (3 classes), la classe neutre a systématiquement le F1
le plus bas. Elle nécessite de détecter une \textit{absence} de sentiment
polarisé, ce qui est plus difficile que d'identifier une émotion forte.
MARBERTv2 est le premier modèle à franchir 0{,}75 sur cette classe.

\subsection{Le coût d'entraînement est négligeable}

Les quatre fine-tunings ont été réalisés gratuitement sur Kaggle en moins
de 12 minutes chacun. Le fine-tuning de DarijaBERT-arabizi (3,4 min) produit
le meilleur score du benchmark — illustrant qu'un bon alignement entre données
de pré-entraînement et données de fine-tuning peut accélérer considérablement
la convergence.

\subsection{Recommandations par scénario}

\begin{center}
\begin{tabular}{lll}
\toprule
\textbf{Scénario} & \textbf{Modèle recommandé} & \textbf{Score} \\
\midrule
Darija marocain (écriture arabe) & MARBERTv2 (fine-tuné) & 0{,}844 \\
Arabizi marocain                 & DarijaBERT-arabizi (fine-tuné) & 0{,}983 \\
Arabe standard (MSA)             & camelbert-da (prêt à l'emploi) & 0{,}924 \\
Français                         & distilcamembert (prêt à l'emploi) & 0{,}949 \\
Modèle unique toutes langues     & xlm-t (acceptable, pas optimal) & 0{,}762--0{,}813 \\
\bottomrule
\end{tabular}
\end{center}

% ============================================================
\section{Prochaines étapes}
% ============================================================

\begin{enumerate}
  \item \textbf{Étape 5 — Atlas-Chat (LLM zéro-shot)} :
    Évaluation du LLM multilingue
    \texttt{MBZUAI-Paris/Atlas-Chat-2B} (modèle ouvert spécialisé arabe
    maghrébin) en mode zéro-shot sur les jeux darija et arabizi. Quantisation
    4 bits (BitsAndBytes) sur GPU T4. Objectif : comparer le coût LLM
    vs le F1 des encodeurs fine-tunés. \textit{Accès conditionnel
    (modèle Gemma, license à confirmer).}

  \item \textbf{Étape 6 — Corpus domaine réel} :
    Constitution d'un jeu de test annoté manuellement (~200 utterances) à
    partir de fichiers SRT de la HACA. Les règles d'annotation seront
    définies avec l'encadrante. Ce jeu permettra de mesurer l'écart entre
    performances sur données académiques et performances sur données réelles.

  \item \textbf{Étape 7 — Rapport final et pipeline SRT} :
    \begin{itemize}
      \item Pipeline SRT complet : parsage → segmentation en utterances →
        routage linguistique (heuristique script) → modèle par langue → résultat
      \item Synthèse état de l'art en TALN multilingue marocain
      \item Rapport final avec recommandations pour intégration dans HMS/HACA Bridges
    \end{itemize}
\end{enumerate}

% ============================================================
\section{Conclusion}
% ============================================================

En quatre étapes, le benchmark a atteint ses objectifs intermédiaires :
jeux de test reproductibles, harness d'évaluation unifié, et identification
des meilleurs modèles pour trois des quatre langues cibles (arabizi, français,
MSA). Pour le darija marocain, le modèle MARBERTv2 fine-tuné dépasse le seuil
opérationnel de 0{,}75 avec un score de \textbf{0{,}844}.

Le résultat le plus marquant est \textbf{DarijaBERT-arabizi à 0{,}983 macro-F1}
sur l'arabizi marocain — une langue pour laquelle aucun modèle dédié
n'existait auparavant. Ce résultat confirme que le fine-tuning de modèles
pré-entraînés sur des données dialectales marocaines constitue la voie
la plus efficace pour des performances opérationnelles.

Les étapes restantes (Atlas-Chat, corpus domaine réel, rapport final)
permettront de valider ces résultats sur des contenus audiovisuels réels
et de formuler des recommandations d'intégration concrètes pour la DSI.

\vspace{1cm}
\begin{center}
\rule{0.4\linewidth}{0.5pt}\\[4pt]
\textit{Marwane ElBaraka — Stage HACA / DSI — Juin 2026}
\end{center}

\end{document}
